\documentclass{ximera}
%herschreven 23 september 2026
\title{Elektrische Veldlijnen}
\author{Daan Lipkens}

\begin{document}
\begin{abstract}
    Het visualiseren van elektrische velden aan de hand van veldlijnen.
\end{abstract}
\maketitle

\section{Veldlijnen tekenen}

Om de wiskundige vectoren van een elektrisch veld inzichtelijk te maken, gebruiken we een grafische voorstelling: \textbf{elektrische veldlijnen} (een concept dat eveneens bedacht werd door Faraday).

Deze veldlijnen geven ons snel een kwalitatief beeld van hoe het veld zich in de ruimte gedraagt en volgen enkele strikte regels:

\begin{proposition}[foldable=true,title={Regels voor elektrische veldlijnen}]\nl
    \begin{enumerate}
        \item De elektrische veldvector $\vec{E}$ is in elk punt volgt de \textbf{raaklijn} aan de veldlijn in dat punt. De pijl op de veldlijn geeft de zin van het veld aan (de richting waarin een positieve testlading weggeduwd zou worden).
        \item De \textbf{dichtheid van de lijnen} (het aantal lijnen per oppervlakte-eenheid) is evenredig met de grootte van de veldsterkte. Veldlijnen die dicht bij elkaar liggen, duiden op een sterk elektrisch veld. Lijnen die ver uit elkaar liggen, duiden op een zwak veld.
        \item Veldlijnen \textbf{vertrekken altijd vanuit een positieve lading} en \textbf{eindigen op een negatieve lading} (of in het oneindige als er geen tegengestelde lading in de buurt is).
        \item Het aantal veldlijnen dat uit een positieve lading vertrekt, of op een negatieve lading toekomt, is \textbf{evenredig met de grootte van die lading}.
        \item Twee elektrische veldlijnen kunnen \textbf{nooit} elkaar snijden!
    \end{enumerate}
\end{proposition}

\begin{remark}[foldable=true]
    Veldlijnen zijn geen echte fysieke objecten in de ruimte, maar slechts een hulpmiddel om het veld te visualiseren. Bovendien stellen ze niet de baan voor die een deeltje zal volgen (behalve in specifieke gevallen zoals beweging op een rechte lijn).
\end{remark}

\begin{remark}[expandable=true, title={Verschil van tekenen tussen vector en veldlijn}]\nl
    Een veldvector (bv. $\vec{E}$) teken je als een rechtelijn met de pijlpunt op het einde. Een veldlijn (bv. $E$) kan een kromme lijn zijn en de pijl staat altijd in het midden.
\end{remark}

\end{document}